% =====================================================================
% m4_verification.tex -- Member 4 (Verification & Memory-Access Traces)
%   * Literature Review, subsection D (trace-driven simulation & cache
%     verification)
%   * Proposed Solution, subsection D (verification methodology & trace
%     design)
%
% For M5: % =====================================================================
% m4_verification.tex -- Member 4 (Verification & Memory-Access Traces)
%   * Literature Review, subsection D (trace-driven simulation & cache
%     verification)
%   * Proposed Solution, subsection D (verification methodology & trace
%     design)
%
% For M5: % =====================================================================
% m4_verification.tex -- Member 4 (Verification & Memory-Access Traces)
%   * Literature Review, subsection D (trace-driven simulation & cache
%     verification)
%   * Proposed Solution, subsection D (verification methodology & trace
%     design)
%
% For M5: % =====================================================================
% m4_verification.tex -- Member 4 (Verification & Memory-Access Traces)
%   * Literature Review, subsection D (trace-driven simulation & cache
%     verification)
%   * Proposed Solution, subsection D (verification methodology & trace
%     design)
%
% For M5: \input{sections/m4_verification} into the IEEEtran report.
% Citation keys used: smith1982cache, hennessy2017quantitative,
%   hill1989associativity, patterson2020organization
% (bibitems at the bottom -- merge into the shared reference list).
% =====================================================================


% --------------- BEGIN LITERATURE REVIEW D (Section 3) ----------------
\subsection{Trace-Driven Simulation and Cache Verification}
\label{sec:m4-litrev}

Evaluating a cache design before it is fabricated requires a way to feed
realistic memory-access patterns into the model and measure the resulting
behaviour.  \emph{Trace-driven simulation} is the classical method: a
sequence of addresses (a \emph{trace}) is recorded or synthesised and then
replayed through a cache simulator.  Smith's landmark survey on cache
memories~\cite{smith1982cache} demonstrated that trace-driven simulation
captures the interplay of block size, associativity and replacement policy
far more faithfully than analytical models, because it preserves the
temporal ordering of references that real programs exhibit.  The same survey
established the two locality principles---temporal and spatial---that
underpin every cache optimisation and that must therefore be represented in
any credible trace set.

Different access patterns expose different aspects of cache behaviour.
Hennessy and Patterson~\cite{hennessy2017quantitative} classify misses into
three categories: \emph{compulsory} (first access to a block),
\emph{capacity} (the working set exceeds the cache size) and
\emph{conflict} (two blocks map to the same index in a direct-mapped or
set-associative cache).  A sequential scan through memory triggers one
compulsory miss per block followed by spatial-locality hits for the
remaining words.  A repeated loop over a small data set converts all
compulsory misses into temporal-locality hits on the second pass.  A random
trace with a large address space stresses capacity, while a carefully
constructed \emph{conflict} trace---in which successive addresses are
separated by exactly the cache size---forces eviction on every access and
drives the hit rate to zero.  Together, these four patterns exercise the
entire miss taxonomy and reveal the strengths and weaknesses of a
particular cache configuration.

Hill and Smith~\cite{hill1989associativity} used trace-driven simulation to
compare hit rates across associativities and cache sizes, introducing the
concept of \emph{inclusion properties} and showing that conflict misses
dominate in small direct-mapped caches.  Their methodology---replaying
large address traces through a parameterised simulator and tabulating hit
rates---is the template adopted in this project.

Verification of the register-transfer-level (RTL) implementation adds a
second dimension.  An RTL simulation (e.g.\ in Icarus Verilog or
ModelSim) produces cycle-accurate waveforms that can be inspected visually
in a Value Change Dump (VCD) viewer and compared automatically against a
\emph{reference model}.  The reference model is typically a software cache
simulator written in a high-level language such as Python.  It accepts the
same trace, simulates the same cache parameters (size, block size, mapping
policy) and records the expected hit/miss outcome for every access.  If the
RTL and the reference model agree on every outcome, the design is
considered functionally correct with respect to that
trace~\cite{patterson2020organization}.  Using multiple, structurally
diverse traces increases confidence that corner cases---such as cold
start-up, back-to-back misses and index conflicts---are covered.
% ---------------- END LITERATURE REVIEW D (Section 3) -----------------


% --------------- BEGIN PROPOSED SOLUTION D (Section 4) ----------------
\subsection{Verification Methodology and Trace Design}
\label{sec:m4-solution}

\subsubsection{Trace generation}
Four trace types are generated by a Python script
(\texttt{generate\_traces.py}) to target each miss category and to produce
a spread of hit rates.  Table~\ref{tab:traces} summarises the patterns and
their expected behaviour for the default 1\,KB direct-mapped cache with
16-byte (4-word) blocks.

\begin{table}[t]
\centering
\caption{Trace types and expected hit rates (1\,KB cache, 16\,B blocks)}
\label{tab:traces}
\begin{tabular}{llcc}
\hline
\textbf{Trace} & \textbf{Pattern} & \textbf{Expected} & \textbf{Dominant}\\
               &                  & \textbf{Hit Rate}  & \textbf{Miss Type}\\
\hline
Sequential  & Linear word scan       & $\approx 75\%$ & Compulsory \\
Repeated    & Small loop, 2+ passes  & $\approx 100\%$\textsuperscript{*} & Compulsory \\
Random      & Uniform 32-bit range   & $\approx 0$--$5\%$ & Capacity \\
Conflict    & Stride = cache size    & $\approx 0\%$ & Conflict \\
\hline
\multicolumn{4}{l}{\footnotesize\textsuperscript{*}After cold (first-pass) misses.}
\end{tabular}
\end{table}

\begin{enumerate}
  \item \textbf{Sequential trace.}  Addresses increase by 4 bytes (one
        word) on every access.  Because each 16-byte block holds four
        words, the first access to a block is a compulsory miss and the
        next three are spatial-locality hits, giving a steady-state hit
        rate of $3/4 = 75\%$.

  \item \textbf{Repeated-loop trace.}  A small set of addresses (fitting
        within distinct cache lines) is traversed repeatedly.  The first
        pass incurs compulsory misses; every subsequent pass hits, so the
        overall hit rate approaches $100\%$ as the number of iterations
        grows.

  \item \textbf{Random trace.}  Addresses are drawn uniformly from a range
        much larger than the cache.  With 64 lines and a 32-bit address
        space, the probability that a random access hits a previously
        loaded and non-evicted line is very low, producing a hit rate in
        the $0$--$5\%$ range.

  \item \textbf{Conflict (thrashing) trace.}  Two or more addresses whose
        index bits are identical (e.g.\ addresses separated by
        $1024 = 2^{10}$ bytes, i.e.\ exactly one cache size apart) are
        alternated.  Each access evicts the line loaded by the previous
        one, so every access after the first is a conflict miss and the
        hit rate is $\approx 0\%$.
\end{enumerate}

Each trace is written in hexadecimal, one 32-bit word-aligned address per
line, and saved in the \texttt{traces/} directory.  The testbench loads a
trace with \texttt{\$readmemh}.

\subsubsection{Testbench architecture}
The Verilog testbench (\texttt{tb\_cache\_controller.v}) instantiates the
cache controller, a behavioural main-memory model and a clock generator.
Its operation proceeds as follows:
\begin{enumerate}
  \item Load a trace file into a register array using \texttt{\$readmemh}.
  \item For each address in the trace, assert a read request and wait for
        the controller to signal \texttt{ready}.
  \item On each completed access, increment either the
        \texttt{hit\_count} or the \texttt{miss\_count} counter based on
        the controller's \texttt{hit} output.
  \item After the last address, compute and display the hit rate:
        $\text{hit rate} = \texttt{hit\_count} \;/\;
        (\texttt{hit\_count} + \texttt{miss\_count})$.
  \item Dump all signal transitions to a VCD file
        (\texttt{\$dumpfile}/\texttt{\$dumpvars}) for waveform inspection.
\end{enumerate}

A \texttt{\$display}-based report is printed at the end of simulation,
showing total accesses, hits, misses and the computed hit rate.  This
allows automated regression: a shell script can run all four traces and
\texttt{grep} the hit rates from the simulation log.

\subsubsection{Python reference model}
A bit-accurate Python reference model (\texttt{cache\_model.py}) mirrors
the RTL cache parameters: cache size, block size, address width and direct
mapping.  It maintains a Python dictionary keyed by index, storing the
valid bit and tag for each line.  For every address in the trace it
performs the same decomposition as the hardware---extracting the tag, index
and offset---and records a hit or miss.  After processing the full trace
it prints the hit rate and a per-access hit/miss log.

Because the model is written in fewer than 80 lines of Python, it is easy
to audit and trust as a \emph{golden reference}.  Discrepancies between
the model and the RTL point to bugs in the Verilog, not in the model.

\subsubsection{End-to-end verification flow}
The complete verification flow is:
\begin{enumerate}
  \item \textbf{Generate} -- run \texttt{generate\_traces.py} to produce
        the four \texttt{.hex} trace files.
  \item \textbf{Simulate} -- compile and run the RTL testbench with Icarus
        Verilog (\texttt{iverilog} + \texttt{vvp}), producing a simulation
        log and a VCD waveform file for each trace.
  \item \textbf{Reference} -- run \texttt{cache\_model.py} on the same
        trace files to obtain golden hit/miss counts.
  \item \textbf{Compare} -- a wrapper script (or manual inspection)
        checks that the RTL hit and miss counts match the reference model
        exactly for every trace.
  \item \textbf{Inspect} -- open the VCD files in GTKWave to visually
        confirm FSM transitions, memory stalls and hit/miss signals on
        selected corner-case accesses.
\end{enumerate}

If all four traces produce matching counts and the waveforms show correct
state-machine behaviour, the cache controller is deemed verified for the
default configuration.  The same flow can be repeated after changing the
cache parameters to verify alternative configurations during Phase~II.
% ---------------- END PROPOSED SOLUTION D (Section 4) -----------------


% -------------- REFERENCES (merge into the shared list) ---------------
% \begin{thebibliography}{9}
% \bibitem{smith1982cache} A.~J. Smith, ``Cache memories,'' \emph{ACM
%   Computing Surveys}, vol.~14, no.~3, pp.~473--530, Sep. 1982.
% \bibitem{hennessy2017quantitative} J.~L. Hennessy and D.~A. Patterson,
%   \emph{Computer Architecture: A Quantitative Approach}, 6th~ed.\
%   Cambridge, MA, USA: Morgan Kaufmann, 2017.
% \bibitem{hill1989associativity} M.~D. Hill and A.~J. Smith, ``Evaluating
%   associativity in CPU caches,'' \emph{IEEE Trans. Computers}, vol.~38,
%   no.~12, pp.~1612--1630, Dec. 1989.
% \bibitem{patterson2020organization} D.~A. Patterson and J.~L. Hennessy,
%   \emph{Computer Organization and Design RISC-V Edition}, 2nd~ed.\
%   Morgan Kaufmann, 2020.
% \end{thebibliography}
 into the IEEEtran report.
% Citation keys used: smith1982cache, hennessy2017quantitative,
%   hill1989associativity, patterson2020organization
% (bibitems at the bottom -- merge into the shared reference list).
% =====================================================================


% --------------- BEGIN LITERATURE REVIEW D (Section 3) ----------------
\subsection{Trace-Driven Simulation and Cache Verification}
\label{sec:m4-litrev}

Evaluating a cache design before it is fabricated requires a way to feed
realistic memory-access patterns into the model and measure the resulting
behaviour.  \emph{Trace-driven simulation} is the classical method: a
sequence of addresses (a \emph{trace}) is recorded or synthesised and then
replayed through a cache simulator.  Smith's landmark survey on cache
memories~\cite{smith1982cache} demonstrated that trace-driven simulation
captures the interplay of block size, associativity and replacement policy
far more faithfully than analytical models, because it preserves the
temporal ordering of references that real programs exhibit.  The same survey
established the two locality principles---temporal and spatial---that
underpin every cache optimisation and that must therefore be represented in
any credible trace set.

Different access patterns expose different aspects of cache behaviour.
Hennessy and Patterson~\cite{hennessy2017quantitative} classify misses into
three categories: \emph{compulsory} (first access to a block),
\emph{capacity} (the working set exceeds the cache size) and
\emph{conflict} (two blocks map to the same index in a direct-mapped or
set-associative cache).  A sequential scan through memory triggers one
compulsory miss per block followed by spatial-locality hits for the
remaining words.  A repeated loop over a small data set converts all
compulsory misses into temporal-locality hits on the second pass.  A random
trace with a large address space stresses capacity, while a carefully
constructed \emph{conflict} trace---in which successive addresses are
separated by exactly the cache size---forces eviction on every access and
drives the hit rate to zero.  Together, these four patterns exercise the
entire miss taxonomy and reveal the strengths and weaknesses of a
particular cache configuration.

Hill and Smith~\cite{hill1989associativity} used trace-driven simulation to
compare hit rates across associativities and cache sizes, introducing the
concept of \emph{inclusion properties} and showing that conflict misses
dominate in small direct-mapped caches.  Their methodology---replaying
large address traces through a parameterised simulator and tabulating hit
rates---is the template adopted in this project.

Verification of the register-transfer-level (RTL) implementation adds a
second dimension.  An RTL simulation (e.g.\ in Icarus Verilog or
ModelSim) produces cycle-accurate waveforms that can be inspected visually
in a Value Change Dump (VCD) viewer and compared automatically against a
\emph{reference model}.  The reference model is typically a software cache
simulator written in a high-level language such as Python.  It accepts the
same trace, simulates the same cache parameters (size, block size, mapping
policy) and records the expected hit/miss outcome for every access.  If the
RTL and the reference model agree on every outcome, the design is
considered functionally correct with respect to that
trace~\cite{patterson2020organization}.  Using multiple, structurally
diverse traces increases confidence that corner cases---such as cold
start-up, back-to-back misses and index conflicts---are covered.
% ---------------- END LITERATURE REVIEW D (Section 3) -----------------


% --------------- BEGIN PROPOSED SOLUTION D (Section 4) ----------------
\subsection{Verification Methodology and Trace Design}
\label{sec:m4-solution}

\subsubsection{Trace generation}
Four trace types are generated by a Python script
(\texttt{generate\_traces.py}) to target each miss category and to produce
a spread of hit rates.  Table~\ref{tab:traces} summarises the patterns and
their expected behaviour for the default 1\,KB direct-mapped cache with
16-byte (4-word) blocks.

\begin{table}[t]
\centering
\caption{Trace types and expected hit rates (1\,KB cache, 16\,B blocks)}
\label{tab:traces}
\begin{tabular}{llcc}
\hline
\textbf{Trace} & \textbf{Pattern} & \textbf{Expected} & \textbf{Dominant}\\
               &                  & \textbf{Hit Rate}  & \textbf{Miss Type}\\
\hline
Sequential  & Linear word scan       & $\approx 75\%$ & Compulsory \\
Repeated    & Small loop, 2+ passes  & $\approx 100\%$\textsuperscript{*} & Compulsory \\
Random      & Uniform 32-bit range   & $\approx 0$--$5\%$ & Capacity \\
Conflict    & Stride = cache size    & $\approx 0\%$ & Conflict \\
\hline
\multicolumn{4}{l}{\footnotesize\textsuperscript{*}After cold (first-pass) misses.}
\end{tabular}
\end{table}

\begin{enumerate}
  \item \textbf{Sequential trace.}  Addresses increase by 4 bytes (one
        word) on every access.  Because each 16-byte block holds four
        words, the first access to a block is a compulsory miss and the
        next three are spatial-locality hits, giving a steady-state hit
        rate of $3/4 = 75\%$.

  \item \textbf{Repeated-loop trace.}  A small set of addresses (fitting
        within distinct cache lines) is traversed repeatedly.  The first
        pass incurs compulsory misses; every subsequent pass hits, so the
        overall hit rate approaches $100\%$ as the number of iterations
        grows.

  \item \textbf{Random trace.}  Addresses are drawn uniformly from a range
        much larger than the cache.  With 64 lines and a 32-bit address
        space, the probability that a random access hits a previously
        loaded and non-evicted line is very low, producing a hit rate in
        the $0$--$5\%$ range.

  \item \textbf{Conflict (thrashing) trace.}  Two or more addresses whose
        index bits are identical (e.g.\ addresses separated by
        $1024 = 2^{10}$ bytes, i.e.\ exactly one cache size apart) are
        alternated.  Each access evicts the line loaded by the previous
        one, so every access after the first is a conflict miss and the
        hit rate is $\approx 0\%$.
\end{enumerate}

Each trace is written in hexadecimal, one 32-bit word-aligned address per
line, and saved in the \texttt{traces/} directory.  The testbench loads a
trace with \texttt{\$readmemh}.

\subsubsection{Testbench architecture}
The Verilog testbench (\texttt{tb\_cache\_controller.v}) instantiates the
cache controller, a behavioural main-memory model and a clock generator.
Its operation proceeds as follows:
\begin{enumerate}
  \item Load a trace file into a register array using \texttt{\$readmemh}.
  \item For each address in the trace, assert a read request and wait for
        the controller to signal \texttt{ready}.
  \item On each completed access, increment either the
        \texttt{hit\_count} or the \texttt{miss\_count} counter based on
        the controller's \texttt{hit} output.
  \item After the last address, compute and display the hit rate:
        $\text{hit rate} = \texttt{hit\_count} \;/\;
        (\texttt{hit\_count} + \texttt{miss\_count})$.
  \item Dump all signal transitions to a VCD file
        (\texttt{\$dumpfile}/\texttt{\$dumpvars}) for waveform inspection.
\end{enumerate}

A \texttt{\$display}-based report is printed at the end of simulation,
showing total accesses, hits, misses and the computed hit rate.  This
allows automated regression: a shell script can run all four traces and
\texttt{grep} the hit rates from the simulation log.

\subsubsection{Python reference model}
A bit-accurate Python reference model (\texttt{cache\_model.py}) mirrors
the RTL cache parameters: cache size, block size, address width and direct
mapping.  It maintains a Python dictionary keyed by index, storing the
valid bit and tag for each line.  For every address in the trace it
performs the same decomposition as the hardware---extracting the tag, index
and offset---and records a hit or miss.  After processing the full trace
it prints the hit rate and a per-access hit/miss log.

Because the model is written in fewer than 80 lines of Python, it is easy
to audit and trust as a \emph{golden reference}.  Discrepancies between
the model and the RTL point to bugs in the Verilog, not in the model.

\subsubsection{End-to-end verification flow}
The complete verification flow is:
\begin{enumerate}
  \item \textbf{Generate} -- run \texttt{generate\_traces.py} to produce
        the four \texttt{.hex} trace files.
  \item \textbf{Simulate} -- compile and run the RTL testbench with Icarus
        Verilog (\texttt{iverilog} + \texttt{vvp}), producing a simulation
        log and a VCD waveform file for each trace.
  \item \textbf{Reference} -- run \texttt{cache\_model.py} on the same
        trace files to obtain golden hit/miss counts.
  \item \textbf{Compare} -- a wrapper script (or manual inspection)
        checks that the RTL hit and miss counts match the reference model
        exactly for every trace.
  \item \textbf{Inspect} -- open the VCD files in GTKWave to visually
        confirm FSM transitions, memory stalls and hit/miss signals on
        selected corner-case accesses.
\end{enumerate}

If all four traces produce matching counts and the waveforms show correct
state-machine behaviour, the cache controller is deemed verified for the
default configuration.  The same flow can be repeated after changing the
cache parameters to verify alternative configurations during Phase~II.
% ---------------- END PROPOSED SOLUTION D (Section 4) -----------------


% -------------- REFERENCES (merge into the shared list) ---------------
% \begin{thebibliography}{9}
% \bibitem{smith1982cache} A.~J. Smith, ``Cache memories,'' \emph{ACM
%   Computing Surveys}, vol.~14, no.~3, pp.~473--530, Sep. 1982.
% \bibitem{hennessy2017quantitative} J.~L. Hennessy and D.~A. Patterson,
%   \emph{Computer Architecture: A Quantitative Approach}, 6th~ed.\
%   Cambridge, MA, USA: Morgan Kaufmann, 2017.
% \bibitem{hill1989associativity} M.~D. Hill and A.~J. Smith, ``Evaluating
%   associativity in CPU caches,'' \emph{IEEE Trans. Computers}, vol.~38,
%   no.~12, pp.~1612--1630, Dec. 1989.
% \bibitem{patterson2020organization} D.~A. Patterson and J.~L. Hennessy,
%   \emph{Computer Organization and Design RISC-V Edition}, 2nd~ed.\
%   Morgan Kaufmann, 2020.
% \end{thebibliography}
 into the IEEEtran report.
% Citation keys used: smith1982cache, hennessy2017quantitative,
%   hill1989associativity, patterson2020organization
% (bibitems at the bottom -- merge into the shared reference list).
% =====================================================================


% --------------- BEGIN LITERATURE REVIEW D (Section 3) ----------------
\subsection{Trace-Driven Simulation and Cache Verification}
\label{sec:m4-litrev}

Evaluating a cache design before it is fabricated requires a way to feed
realistic memory-access patterns into the model and measure the resulting
behaviour.  \emph{Trace-driven simulation} is the classical method: a
sequence of addresses (a \emph{trace}) is recorded or synthesised and then
replayed through a cache simulator.  Smith's landmark survey on cache
memories~\cite{smith1982cache} demonstrated that trace-driven simulation
captures the interplay of block size, associativity and replacement policy
far more faithfully than analytical models, because it preserves the
temporal ordering of references that real programs exhibit.  The same survey
established the two locality principles---temporal and spatial---that
underpin every cache optimisation and that must therefore be represented in
any credible trace set.

Different access patterns expose different aspects of cache behaviour.
Hennessy and Patterson~\cite{hennessy2017quantitative} classify misses into
three categories: \emph{compulsory} (first access to a block),
\emph{capacity} (the working set exceeds the cache size) and
\emph{conflict} (two blocks map to the same index in a direct-mapped or
set-associative cache).  A sequential scan through memory triggers one
compulsory miss per block followed by spatial-locality hits for the
remaining words.  A repeated loop over a small data set converts all
compulsory misses into temporal-locality hits on the second pass.  A random
trace with a large address space stresses capacity, while a carefully
constructed \emph{conflict} trace---in which successive addresses are
separated by exactly the cache size---forces eviction on every access and
drives the hit rate to zero.  Together, these four patterns exercise the
entire miss taxonomy and reveal the strengths and weaknesses of a
particular cache configuration.

Hill and Smith~\cite{hill1989associativity} used trace-driven simulation to
compare hit rates across associativities and cache sizes, introducing the
concept of \emph{inclusion properties} and showing that conflict misses
dominate in small direct-mapped caches.  Their methodology---replaying
large address traces through a parameterised simulator and tabulating hit
rates---is the template adopted in this project.

Verification of the register-transfer-level (RTL) implementation adds a
second dimension.  An RTL simulation (e.g.\ in Icarus Verilog or
ModelSim) produces cycle-accurate waveforms that can be inspected visually
in a Value Change Dump (VCD) viewer and compared automatically against a
\emph{reference model}.  The reference model is typically a software cache
simulator written in a high-level language such as Python.  It accepts the
same trace, simulates the same cache parameters (size, block size, mapping
policy) and records the expected hit/miss outcome for every access.  If the
RTL and the reference model agree on every outcome, the design is
considered functionally correct with respect to that
trace~\cite{patterson2020organization}.  Using multiple, structurally
diverse traces increases confidence that corner cases---such as cold
start-up, back-to-back misses and index conflicts---are covered.
% ---------------- END LITERATURE REVIEW D (Section 3) -----------------


% --------------- BEGIN PROPOSED SOLUTION D (Section 4) ----------------
\subsection{Verification Methodology and Trace Design}
\label{sec:m4-solution}

\subsubsection{Trace generation}
Four trace types are generated by a Python script
(\texttt{generate\_traces.py}) to target each miss category and to produce
a spread of hit rates.  Table~\ref{tab:traces} summarises the patterns and
their expected behaviour for the default 1\,KB direct-mapped cache with
16-byte (4-word) blocks.

\begin{table}[t]
\centering
\caption{Trace types and expected hit rates (1\,KB cache, 16\,B blocks)}
\label{tab:traces}
\begin{tabular}{llcc}
\hline
\textbf{Trace} & \textbf{Pattern} & \textbf{Expected} & \textbf{Dominant}\\
               &                  & \textbf{Hit Rate}  & \textbf{Miss Type}\\
\hline
Sequential  & Linear word scan       & $\approx 75\%$ & Compulsory \\
Repeated    & Small loop, 2+ passes  & $\approx 100\%$\textsuperscript{*} & Compulsory \\
Random      & Uniform 32-bit range   & $\approx 0$--$5\%$ & Capacity \\
Conflict    & Stride = cache size    & $\approx 0\%$ & Conflict \\
\hline
\multicolumn{4}{l}{\footnotesize\textsuperscript{*}After cold (first-pass) misses.}
\end{tabular}
\end{table}

\begin{enumerate}
  \item \textbf{Sequential trace.}  Addresses increase by 4 bytes (one
        word) on every access.  Because each 16-byte block holds four
        words, the first access to a block is a compulsory miss and the
        next three are spatial-locality hits, giving a steady-state hit
        rate of $3/4 = 75\%$.

  \item \textbf{Repeated-loop trace.}  A small set of addresses (fitting
        within distinct cache lines) is traversed repeatedly.  The first
        pass incurs compulsory misses; every subsequent pass hits, so the
        overall hit rate approaches $100\%$ as the number of iterations
        grows.

  \item \textbf{Random trace.}  Addresses are drawn uniformly from a range
        much larger than the cache.  With 64 lines and a 32-bit address
        space, the probability that a random access hits a previously
        loaded and non-evicted line is very low, producing a hit rate in
        the $0$--$5\%$ range.

  \item \textbf{Conflict (thrashing) trace.}  Two or more addresses whose
        index bits are identical (e.g.\ addresses separated by
        $1024 = 2^{10}$ bytes, i.e.\ exactly one cache size apart) are
        alternated.  Each access evicts the line loaded by the previous
        one, so every access after the first is a conflict miss and the
        hit rate is $\approx 0\%$.
\end{enumerate}

Each trace is written in hexadecimal, one 32-bit word-aligned address per
line, and saved in the \texttt{traces/} directory.  The testbench loads a
trace with \texttt{\$readmemh}.

\subsubsection{Testbench architecture}
The Verilog testbench (\texttt{tb\_cache\_controller.v}) instantiates the
cache controller, a behavioural main-memory model and a clock generator.
Its operation proceeds as follows:
\begin{enumerate}
  \item Load a trace file into a register array using \texttt{\$readmemh}.
  \item For each address in the trace, assert a read request and wait for
        the controller to signal \texttt{ready}.
  \item On each completed access, increment either the
        \texttt{hit\_count} or the \texttt{miss\_count} counter based on
        the controller's \texttt{hit} output.
  \item After the last address, compute and display the hit rate:
        $\text{hit rate} = \texttt{hit\_count} \;/\;
        (\texttt{hit\_count} + \texttt{miss\_count})$.
  \item Dump all signal transitions to a VCD file
        (\texttt{\$dumpfile}/\texttt{\$dumpvars}) for waveform inspection.
\end{enumerate}

A \texttt{\$display}-based report is printed at the end of simulation,
showing total accesses, hits, misses and the computed hit rate.  This
allows automated regression: a shell script can run all four traces and
\texttt{grep} the hit rates from the simulation log.

\subsubsection{Python reference model}
A bit-accurate Python reference model (\texttt{cache\_model.py}) mirrors
the RTL cache parameters: cache size, block size, address width and direct
mapping.  It maintains a Python dictionary keyed by index, storing the
valid bit and tag for each line.  For every address in the trace it
performs the same decomposition as the hardware---extracting the tag, index
and offset---and records a hit or miss.  After processing the full trace
it prints the hit rate and a per-access hit/miss log.

Because the model is written in fewer than 80 lines of Python, it is easy
to audit and trust as a \emph{golden reference}.  Discrepancies between
the model and the RTL point to bugs in the Verilog, not in the model.

\subsubsection{End-to-end verification flow}
The complete verification flow is:
\begin{enumerate}
  \item \textbf{Generate} -- run \texttt{generate\_traces.py} to produce
        the four \texttt{.hex} trace files.
  \item \textbf{Simulate} -- compile and run the RTL testbench with Icarus
        Verilog (\texttt{iverilog} + \texttt{vvp}), producing a simulation
        log and a VCD waveform file for each trace.
  \item \textbf{Reference} -- run \texttt{cache\_model.py} on the same
        trace files to obtain golden hit/miss counts.
  \item \textbf{Compare} -- a wrapper script (or manual inspection)
        checks that the RTL hit and miss counts match the reference model
        exactly for every trace.
  \item \textbf{Inspect} -- open the VCD files in GTKWave to visually
        confirm FSM transitions, memory stalls and hit/miss signals on
        selected corner-case accesses.
\end{enumerate}

If all four traces produce matching counts and the waveforms show correct
state-machine behaviour, the cache controller is deemed verified for the
default configuration.  The same flow can be repeated after changing the
cache parameters to verify alternative configurations during Phase~II.
% ---------------- END PROPOSED SOLUTION D (Section 4) -----------------


% -------------- REFERENCES (merge into the shared list) ---------------
% \begin{thebibliography}{9}
% \bibitem{smith1982cache} A.~J. Smith, ``Cache memories,'' \emph{ACM
%   Computing Surveys}, vol.~14, no.~3, pp.~473--530, Sep. 1982.
% \bibitem{hennessy2017quantitative} J.~L. Hennessy and D.~A. Patterson,
%   \emph{Computer Architecture: A Quantitative Approach}, 6th~ed.\
%   Cambridge, MA, USA: Morgan Kaufmann, 2017.
% \bibitem{hill1989associativity} M.~D. Hill and A.~J. Smith, ``Evaluating
%   associativity in CPU caches,'' \emph{IEEE Trans. Computers}, vol.~38,
%   no.~12, pp.~1612--1630, Dec. 1989.
% \bibitem{patterson2020organization} D.~A. Patterson and J.~L. Hennessy,
%   \emph{Computer Organization and Design RISC-V Edition}, 2nd~ed.\
%   Morgan Kaufmann, 2020.
% \end{thebibliography}
 into the IEEEtran report.
% Citation keys used: smith1982cache, hennessy2017quantitative,
%   hill1989associativity, patterson2020organization
% (bibitems at the bottom -- merge into the shared reference list).
% =====================================================================


% --------------- BEGIN LITERATURE REVIEW D (Section 3) ----------------
\subsection{Trace-Driven Simulation and Cache Verification}
\label{sec:m4-litrev}

Evaluating a cache design before it is fabricated requires a way to feed
realistic memory-access patterns into the model and measure the resulting
behaviour.  \emph{Trace-driven simulation} is the classical method: a
sequence of addresses (a \emph{trace}) is recorded or synthesised and then
replayed through a cache simulator.  Smith's landmark survey on cache
memories~\cite{smith1982cache} demonstrated that trace-driven simulation
captures the interplay of block size, associativity and replacement policy
far more faithfully than analytical models, because it preserves the
temporal ordering of references that real programs exhibit.  The same survey
established the two locality principles---temporal and spatial---that
underpin every cache optimisation and that must therefore be represented in
any credible trace set.

Different access patterns expose different aspects of cache behaviour.
Hennessy and Patterson~\cite{hennessy2017quantitative} classify misses into
three categories: \emph{compulsory} (first access to a block),
\emph{capacity} (the working set exceeds the cache size) and
\emph{conflict} (two blocks map to the same index in a direct-mapped or
set-associative cache).  A sequential scan through memory triggers one
compulsory miss per block followed by spatial-locality hits for the
remaining words.  A repeated loop over a small data set converts all
compulsory misses into temporal-locality hits on the second pass.  A random
trace with a large address space stresses capacity, while a carefully
constructed \emph{conflict} trace---in which successive addresses are
separated by exactly the cache size---forces eviction on every access and
drives the hit rate to zero.  Together, these four patterns exercise the
entire miss taxonomy and reveal the strengths and weaknesses of a
particular cache configuration.

Hill and Smith~\cite{hill1989associativity} used trace-driven simulation to
compare hit rates across associativities and cache sizes, introducing the
concept of \emph{inclusion properties} and showing that conflict misses
dominate in small direct-mapped caches.  Their methodology---replaying
large address traces through a parameterised simulator and tabulating hit
rates---is the template adopted in this project.

Verification of the register-transfer-level (RTL) implementation adds a
second dimension.  An RTL simulation (e.g.\ in Icarus Verilog or
ModelSim) produces cycle-accurate waveforms that can be inspected visually
in a Value Change Dump (VCD) viewer and compared automatically against a
\emph{reference model}.  The reference model is typically a software cache
simulator written in a high-level language such as Python.  It accepts the
same trace, simulates the same cache parameters (size, block size, mapping
policy) and records the expected hit/miss outcome for every access.  If the
RTL and the reference model agree on every outcome, the design is
considered functionally correct with respect to that
trace~\cite{patterson2020organization}.  Using multiple, structurally
diverse traces increases confidence that corner cases---such as cold
start-up, back-to-back misses and index conflicts---are covered.
% ---------------- END LITERATURE REVIEW D (Section 3) -----------------


% --------------- BEGIN PROPOSED SOLUTION D (Section 4) ----------------
\subsection{Verification Methodology and Trace Design}
\label{sec:m4-solution}

\subsubsection{Trace generation}
Four trace types are generated by a Python script
(\texttt{generate\_traces.py}) to target each miss category and to produce
a spread of hit rates.  Table~\ref{tab:traces} summarises the patterns and
their expected behaviour for the default 1\,KB direct-mapped cache with
16-byte (4-word) blocks.

\begin{table}[t]
\centering
\caption{Trace types and expected hit rates (1\,KB cache, 16\,B blocks)}
\label{tab:traces}
\begin{tabular}{llcc}
\hline
\textbf{Trace} & \textbf{Pattern} & \textbf{Expected} & \textbf{Dominant}\\
               &                  & \textbf{Hit Rate}  & \textbf{Miss Type}\\
\hline
Sequential  & Linear word scan       & $\approx 75\%$ & Compulsory \\
Repeated    & Small loop, 2+ passes  & $\approx 100\%$\textsuperscript{*} & Compulsory \\
Random      & Uniform 32-bit range   & $\approx 0$--$5\%$ & Capacity \\
Conflict    & Stride = cache size    & $\approx 0\%$ & Conflict \\
\hline
\multicolumn{4}{l}{\footnotesize\textsuperscript{*}After cold (first-pass) misses.}
\end{tabular}
\end{table}

\begin{enumerate}
  \item \textbf{Sequential trace.}  Addresses increase by 4 bytes (one
        word) on every access.  Because each 16-byte block holds four
        words, the first access to a block is a compulsory miss and the
        next three are spatial-locality hits, giving a steady-state hit
        rate of $3/4 = 75\%$.

  \item \textbf{Repeated-loop trace.}  A small set of addresses (fitting
        within distinct cache lines) is traversed repeatedly.  The first
        pass incurs compulsory misses; every subsequent pass hits, so the
        overall hit rate approaches $100\%$ as the number of iterations
        grows.

  \item \textbf{Random trace.}  Addresses are drawn uniformly from a range
        much larger than the cache.  With 64 lines and a 32-bit address
        space, the probability that a random access hits a previously
        loaded and non-evicted line is very low, producing a hit rate in
        the $0$--$5\%$ range.

  \item \textbf{Conflict (thrashing) trace.}  Two or more addresses whose
        index bits are identical (e.g.\ addresses separated by
        $1024 = 2^{10}$ bytes, i.e.\ exactly one cache size apart) are
        alternated.  Each access evicts the line loaded by the previous
        one, so every access after the first is a conflict miss and the
        hit rate is $\approx 0\%$.
\end{enumerate}

Each trace is written in hexadecimal, one 32-bit word-aligned address per
line, and saved in the \texttt{traces/} directory.  The testbench loads a
trace with \texttt{\$readmemh}.

\subsubsection{Testbench architecture}
The Verilog testbench (\texttt{tb\_cache\_controller.v}) instantiates the
cache controller, a behavioural main-memory model and a clock generator.
Its operation proceeds as follows:
\begin{enumerate}
  \item Load a trace file into a register array using \texttt{\$readmemh}.
  \item For each address in the trace, assert a read request and wait for
        the controller to signal \texttt{ready}.
  \item On each completed access, increment either the
        \texttt{hit\_count} or the \texttt{miss\_count} counter based on
        the controller's \texttt{hit} output.
  \item After the last address, compute and display the hit rate:
        $\text{hit rate} = \texttt{hit\_count} \;/\;
        (\texttt{hit\_count} + \texttt{miss\_count})$.
  \item Dump all signal transitions to a VCD file
        (\texttt{\$dumpfile}/\texttt{\$dumpvars}) for waveform inspection.
\end{enumerate}

A \texttt{\$display}-based report is printed at the end of simulation,
showing total accesses, hits, misses and the computed hit rate.  This
allows automated regression: a shell script can run all four traces and
\texttt{grep} the hit rates from the simulation log.

\subsubsection{Python reference model}
A bit-accurate Python reference model (\texttt{cache\_model.py}) mirrors
the RTL cache parameters: cache size, block size, address width and direct
mapping.  It maintains a Python dictionary keyed by index, storing the
valid bit and tag for each line.  For every address in the trace it
performs the same decomposition as the hardware---extracting the tag, index
and offset---and records a hit or miss.  After processing the full trace
it prints the hit rate and a per-access hit/miss log.

Because the model is written in fewer than 80 lines of Python, it is easy
to audit and trust as a \emph{golden reference}.  Discrepancies between
the model and the RTL point to bugs in the Verilog, not in the model.

\subsubsection{End-to-end verification flow}
The complete verification flow is:
\begin{enumerate}
  \item \textbf{Generate} -- run \texttt{generate\_traces.py} to produce
        the four \texttt{.hex} trace files.
  \item \textbf{Simulate} -- compile and run the RTL testbench with Icarus
        Verilog (\texttt{iverilog} + \texttt{vvp}), producing a simulation
        log and a VCD waveform file for each trace.
  \item \textbf{Reference} -- run \texttt{cache\_model.py} on the same
        trace files to obtain golden hit/miss counts.
  \item \textbf{Compare} -- a wrapper script (or manual inspection)
        checks that the RTL hit and miss counts match the reference model
        exactly for every trace.
  \item \textbf{Inspect} -- open the VCD files in GTKWave to visually
        confirm FSM transitions, memory stalls and hit/miss signals on
        selected corner-case accesses.
\end{enumerate}

If all four traces produce matching counts and the waveforms show correct
state-machine behaviour, the cache controller is deemed verified for the
default configuration.  The same flow can be repeated after changing the
cache parameters to verify alternative configurations during Phase~II.
% ---------------- END PROPOSED SOLUTION D (Section 4) -----------------


% -------------- REFERENCES (merge into the shared list) ---------------
% \begin{thebibliography}{9}
% \bibitem{smith1982cache} A.~J. Smith, ``Cache memories,'' \emph{ACM
%   Computing Surveys}, vol.~14, no.~3, pp.~473--530, Sep. 1982.
% \bibitem{hennessy2017quantitative} J.~L. Hennessy and D.~A. Patterson,
%   \emph{Computer Architecture: A Quantitative Approach}, 6th~ed.\
%   Cambridge, MA, USA: Morgan Kaufmann, 2017.
% \bibitem{hill1989associativity} M.~D. Hill and A.~J. Smith, ``Evaluating
%   associativity in CPU caches,'' \emph{IEEE Trans. Computers}, vol.~38,
%   no.~12, pp.~1612--1630, Dec. 1989.
% \bibitem{patterson2020organization} D.~A. Patterson and J.~L. Hennessy,
%   \emph{Computer Organization and Design RISC-V Edition}, 2nd~ed.\
%   Morgan Kaufmann, 2020.
% \end{thebibliography}
